\section{Bounds and extremal elements}

\import{relation/properties.tex}
\import{order/order.tex}

The following definitions use the reflexive closure of the displayed
relation.  They therefore apply uniformly to strict and non-strict order
relations.

\begin{definition}\label{lower_bound}
    $a$ is a lower bound of $X$ in $A$ with respect to $R$ iff
        $a\in A$ and for all $x\in X$ we have $a=x$ or $a\mathrel{R}x$.
\end{definition}

\begin{definition}\label{upper_bound}
    $a$ is an upper bound of $X$ in $A$ with respect to $R$ iff
        $a\in A$ and for all $x\in X$ we have $x=a$ or $x\mathrel{R}a$.
\end{definition}

\begin{proposition}\label{upper_bound_elim}
    Suppose $a$ is an upper bound of $X$ in $A$ with respect to $R$.
    Suppose $x\in X$.
    Then $x=a$ or $x\mathrel{R}a$.
\end{proposition}
\begin{proof}
    $a\in A$ and for all $y\in X$ we have
        $y=a$ or $y\mathrel{R}a$ by \obv{upper_bound}.
\end{proof}

\begin{definition}\label{least_upper_bound}
    $u$ is a least upper bound of $X$ in $A$ with respect to $R$ iff
        $u$ is an upper bound of $X$ in $A$ with respect to $R$ and
        for every upper bound $v$ of $X$ in $A$ with respect to $R$
            we have $u=v$ or $u\mathrel{R}v$.
\end{definition}

\begin{definition}\label{greatest_lower_bound}
    $u$ is a greatest lower bound of $X$ in $A$ with respect to $R$ iff
        $u$ is a lower bound of $X$ in $A$ with respect to $R$ and
        for every lower bound $v$ of $X$ in $A$ with respect to $R$
            we have $u=v$ or $v\mathrel{R}u$.
\end{definition}

\begin{proposition}\label{least_upper_bound_is_upper_bound}
    Suppose $u$ is a least upper bound of $X$ in $A$ with respect to $R$.
    Then $u$ is an upper bound of $X$ in $A$ with respect to $R$.
\end{proposition}
\begin{proof}
    Follows by \obv{least_upper_bound}.
\end{proof}

\begin{proposition}\label{least_upper_bound_is_least}
    Suppose $u$ is a least upper bound of $X$ in $A$ with respect to $R$ and
        $v$ is an upper bound of $X$ in $A$ with respect to $R$.
    Then $u=v$ or $u\mathrel{R}v$.
\end{proposition}
\begin{proof}
    Follows by \obv{least_upper_bound}.
\end{proof}

\begin{proposition}\label{greatest_lower_bound_is_lower_bound}
    Suppose $u$ is a greatest lower bound of $X$ in $A$ with respect to $R$.
    Then $u$ is a lower bound of $X$ in $A$ with respect to $R$.
\end{proposition}
\begin{proof}
    Follows by \obv{greatest_lower_bound}.
\end{proof}

\begin{proposition}\label{greatest_lower_bound_is_greatest}
    Suppose $u$ is a greatest lower bound of $X$ in $A$ with respect to $R$ and
        $v$ is a lower bound of $X$ in $A$ with respect to $R$.
    Then $u=v$ or $v\mathrel{R}u$.
\end{proposition}
\begin{proof}
    Follows by \obv{greatest_lower_bound}.
\end{proof}

\begin{proposition}\label{least_upper_bound_unique}
    Suppose $R$ is antisymmetric and
        $u$ and $v$ are least upper bounds of $X$ in $A$ with respect to $R$.
    Then $u=v$.
\end{proposition}
\begin{proof}
    Follows by \auto{least_upper_bound,antisymmetric}.
\end{proof}

\begin{proposition}\label{greatest_lower_bound_unique}
    Suppose $R$ is antisymmetric and
        $u$ and $v$ are greatest lower bounds of $X$ in $A$ with respect to $R$.
    Then $u=v$.
\end{proposition}
\begin{proof}
    Follows by \auto{greatest_lower_bound,antisymmetric}.
\end{proof}

\begin{proposition}\label{upper_bound_iff_reflexive_on}
    Suppose $R$ is reflexive on $A$.
    Then $u$ is an upper bound of $X$ in $A$ with respect to $R$ iff
        $u\in A$ and for all $x\in X$ we have $x\mathrel{R}u$.
\end{proposition}
\begin{proof}
    Follows by \auto{upper_bound,reflexive_on}.
\end{proof}

\begin{proposition}\label{lower_bound_iff_reflexive_on}
    Suppose $R$ is reflexive on $A$.
    Then $u$ is a lower bound of $X$ in $A$ with respect to $R$ iff
        $u\in A$ and for all $x\in X$ we have $u\mathrel{R}x$.
\end{proposition}
\begin{proof}
    Follows by \auto{lower_bound,reflexive_on}.
\end{proof}

\begin{proposition}\label{least_upper_bound_iff_order_on}
    Suppose $R$ is an order on $A$.
    Then $u$ is a least upper bound of $X$ in $A$ with respect to $R$ iff
        $u\in A\land
        (\forall x:\Set.\ x\in X\implies x\mathrel{R}u)\land
        (\forall v:\Set.\ (v\in A\land
                (\forall y:\Set.\ y\in X\implies y\mathrel{R}v))
            \implies u\mathrel{R}v)$.
\end{proposition}
\begin{proof}
    $R$ is reflexive on $A$ by definition.
    Follows by
        \auto{least_upper_bound,upper_bound_iff_reflexive_on,
            reflexive_on}.
\end{proof}

\begin{proposition}\label{greatest_lower_bound_iff_order_on}
    Suppose $R$ is an order on $A$.
    Then $u$ is a greatest lower bound of $X$ in $A$ with respect to $R$ iff
        $u\in A\land
        (\forall x:\Set.\ x\in X\implies u\mathrel{R}x)\land
        (\forall v:\Set.\ (v\in A\land
                (\forall y:\Set.\ y\in X\implies v\mathrel{R}y))
            \implies v\mathrel{R}u)$.
\end{proposition}
\begin{proof}
    $R$ is reflexive on $A$ by definition.
    Follows by
        \auto{greatest_lower_bound,lower_bound_iff_reflexive_on,
            reflexive_on}.
\end{proof}

\begin{definition}\label{least_element}
    $a$ is a least element of $X$ with respect to $R$ iff
        $a\in X$ and for all $x\in X$ we have $a=x$ or $a\mathrel{R}x$.
\end{definition}

\begin{definition}\label{greatest_element}
    $a$ is a greatest element of $X$ with respect to $R$ iff
        $a\in X$ and for all $x\in X$ we have $x=a$ or $x\mathrel{R}a$.
\end{definition}

\begin{definition}\label{minimal_element}
    $a$ is a minimal element of $X$ with respect to $R$ iff
        $a\in X$ and for all $x\in X$ such that $x\mathrel{R}a$ we have
            $x=a$.
\end{definition}

\begin{definition}\label{maximal_element}
    $a$ is a maximal element of $X$ with respect to $R$ iff
        $a\in X$ and for all $x\in X$ such that $a\mathrel{R}x$ we have
            $x=a$.
\end{definition}

\begin{proposition}\label{least_iff_lower_bound_in_self}
    $a$ is a least element of $X$ with respect to $R$ iff
        $a\in X$ and $a$ is a lower bound of $X$ in $X$ with respect to $R$.
\end{proposition}
\begin{proof}
    Follows by \auto{least_element,lower_bound}.
\end{proof}

\begin{proposition}\label{greatest_iff_upper_bound_in_self}
    $a$ is a greatest element of $X$ with respect to $R$ iff
        $a\in X$ and $a$ is an upper bound of $X$ in $X$ with respect to $R$.
\end{proposition}
\begin{proof}
    Follows by \auto{greatest_element,upper_bound}.
\end{proof}

\begin{proposition}\label{lower_bound_antitone}
    Suppose $X\subseteq Y$ and
        $a$ is a lower bound of $Y$ in $A$ with respect to $R$.
    Then $a$ is a lower bound of $X$ in $A$ with respect to $R$.
\end{proposition}
\begin{proof}
    Follows by \obv{lower_bound,subseteq}.
\end{proof}

\begin{proposition}\label{upper_bound_antitone}
    Suppose $X\subseteq Y$ and
        $a$ is an upper bound of $Y$ in $A$ with respect to $R$.
    Then $a$ is an upper bound of $X$ in $A$ with respect to $R$.
\end{proposition}
\begin{proof}
    Follows by \obv{upper_bound,subseteq}.
\end{proof}

\begin{proposition}\label{singleton_least_element}
    $a$ is a least element of $\{a\}$ with respect to $R$.
\end{proposition}
\begin{proof}
    Follows by \auto{least_element,singleton_iff}.
\end{proof}

\begin{proposition}\label{singleton_greatest_element}
    $a$ is a greatest element of $\{a\}$ with respect to $R$.
\end{proposition}
\begin{proof}
    Follows by \auto{greatest_element,singleton_iff}.
\end{proof}

\begin{proposition}\label{least_element_is_minimal}
    Suppose $R$ is antisymmetric and
        $a$ is a least element of $X$ with respect to $R$.
    Then $a$ is a minimal element of $X$ with respect to $R$.
\end{proposition}
\begin{proof}
    Follows by \auto{least_element,minimal_element,antisymmetric}.
\end{proof}

\begin{proposition}\label{greatest_element_is_maximal}
    Suppose $R$ is antisymmetric and
        $a$ is a greatest element of $X$ with respect to $R$.
    Then $a$ is a maximal element of $X$ with respect to $R$.
\end{proposition}
\begin{proof}
    Follows by \auto{greatest_element,maximal_element,antisymmetric}.
\end{proof}

\begin{proposition}\label{least_element_unique}
    Suppose $R$ is antisymmetric and
        $a$ and $b$ are least elements of $X$ with respect to $R$.
    Then $a=b$.
\end{proposition}
\begin{proof}
    Follows by \auto{least_element,antisymmetric}.
\end{proof}

\begin{proposition}\label{greatest_element_unique}
    Suppose $R$ is antisymmetric and
        $a$ and $b$ are greatest elements of $X$ with respect to $R$.
    Then $a=b$.
\end{proposition}
\begin{proof}
    Follows by \auto{greatest_element,antisymmetric}.
\end{proof}
